\subsection{THE THEORY OF SETS}

The theory of sets is a theory which contains the relational signs \(=,\) \(\in\) and the substantific sign \(\supset\) (all these signs being of weight 2); in addition to the schemes S1 to S7 given in Chapter I, it contains the scheme S8, which will be introduced in no. 6, and the explicit axioms A1 (no. 3.) A2 (no. 5), A3 ( \(\S\) 2, no. 1), A4 ( \(\S\) 5, no. 1), and A5 (Chapter III, \(\S\) 6, no. 1), These explicit axioms contain no letters; in other words, the theory of sets is a theory without constants.

Since the theory of sets is an equalitarian theory, the results of Chapter I are applicable.

From now on, unless the contrary is expressly stated, we shall always argue in a theory which is stronger (Chapter I, §2, no. 4) than the theory of sets; if the theory is not mentioned, it is to be assumed that the theory of sets is implied. It will be evident in many cases that this hypothesis is superfluous, and the reader should have no difficulty in determining in what theory weaker than the theory of sets the results stated are valid.

If T and U are terms, the assembly \(\in TU\) is a relation (called the relation of membership) which in practice we write in one of the following ways: \(\widetilde{T} \in U\) , \((\widetilde{T}) \in (U)\) , ``T belongs to U'', ``T is an element of U''. The relation ``not \((\widetilde{T} \in U)\) '' is written \(T \notin U\) .

From a ``naive'' point of view, many mathematical entities can be considered as collections or ``sets'' of objects. We do not seek to formalize this notion, and in the formalistic interpretation of what follows, the word ``set'' is to be considered as strictly synonymous with ``term''. In particular, phrases such as ``let X be a set'' are, in principle, quite superfluous, since every letter is a term. Such phrases are introduced only to assist the intuitive interpretation of the text.

